\documentclass[11pt]{article}
\usepackage[letterpaper,margin=1in]{geometry}
\usepackage[T1]{fontenc}
\usepackage{lmodern,microtype}
\usepackage{amsmath,amssymb,amsthm,mathtools}
\usepackage{booktabs,longtable,array}
\usepackage{xcolor}
\usepackage{enumitem}
\usepackage{hyperref}
\hypersetup{colorlinks=true,linkcolor=blue!45!black,urlcolor=blue!55!black,citecolor=blue!45!black,pdftitle={A shared Pell certificate for Con(ZFC)},pdfauthor={Research prepared with ChatGPT for Vladimir Reshetnikov}}
\setlength{\emergencystretch}{3em}
\setlist{itemsep=3pt,topsep=5pt}
\newtheorem{theorem}{Theorem}[section]
\newtheorem{lemma}[theorem]{Lemma}
\newtheorem{proposition}[theorem]{Proposition}
\newtheorem{corollary}[theorem]{Corollary}
\theoremstyle{definition}
\newtheorem{definition}[theorem]{Definition}
\theoremstyle{remark}
\newtheorem{remark}[theorem]{Remark}
\newcommand{\NN}{\mathbb N}
\newcommand{\ZZ}{\mathbb Z}
\newcommand{\Con}{\operatorname{Con}}
\newcommand{\Refut}{\operatorname{Ref}}
\newcommand{\Norm}[3]{#2^2-(#1^2-1)#3^2}
\newcommand{\File}[1]{\nolinkurl{#1}}
\newcommand{\I}{\mathsf I}
\newcommand{\E}{\mathsf E}
\newcommand{\Q}{\mathsf Q}
\newcommand{\J}{\mathsf J}
\newcommand{\Neg}{\mathsf N}
\newcommand{\Ex}{\mathsf{Ex}}
\title{A shared Pell certificate for $\Con(\mathrm{ZFC})$\\[4pt]
\large An explicit quartic in 1,873 unrestricted integer variables}
\author{Research prepared for Vladimir Reshetnikov\\with ChatGPT assistance}
\date{23 October 2026\\\small Repository baseline: \texttt{b594a5e988e986f695923d068ef3d33ab2ada16f}}
\begin{document}
\maketitle
\begin{abstract}
We give a completely materialized integer polynomial $P$ of degree four,
in 1,873 unrestricted integer variables, for which an ordinary proof
formalizable in ZFC establishes
\[
 \mathrm{ZFC}\vdash
 \bigl[\forall\mathbf z\in\ZZ^{1873}\;P(\mathbf z)\ne0\bigr]
 \longleftrightarrow \Con(\mathrm{ZFC}).
\]
The polynomial is a sum of 1,468 explicitly supplied quadratic squares;
its collected expansion has 77,523 nonzero monomials and coefficient
height $1{,}323{,}565{,}948{,}780{,}873{,}728<2^{61}$.
The complete coefficients, deterministic generator, independent exact
checkers, source interfaces and proof notes accompany this article.

The main new result is a simultaneous Pell certificate: forty positive
power graphs share one bootstrap, one auxiliary Pell solution and one
CRT base. Equal exponents share coordinates; three nearby exponents
are recovered by recurrences. A signed index recovery lemma removes
both auxiliary sign guards for each of sixteen seed indices. A single
common output bound replaces forty private bounds. These are
existence-preserving mathematical transformations with explicit
completeness constructions, rather than informal identification of
unrelated existential witnesses.

The fixed logical input is the repository's C222 numerical calculus.
Its established ZFC-refutation calibration is inherited. No K215
simulation or Jones universal polynomial is used. Relative to the latest
incoming direct quartic, the integer-variable count falls from 7,617 to
1,873. This is a reduction within the direct calculus construction;
it is not a claim of a globally smallest consistency equation or a new
minimum among all quartics. No new Lean or encoded object-language ZFC
proof certificate is claimed.
\end{abstract}
\newpage
\tableofcontents
\newpage

\section{The result and its exact scope}

The user-facing equation is the following ordinary polynomial equation:
\begin{equation}\label{eq:final}
 \boxed{P(z_0,\ldots,z_{1872})=
       \sum_{\nu=0}^{1467} Q_\nu(z_0,\ldots,z_{1872})^2=0,\qquad
       (z_0,\ldots,z_{1872})\in\ZZ^{1873}.}
\end{equation}
Every $Q_\nu$ is explicitly listed in \File{equations/residuals.json}.
The collected integer coefficients of $P$ itself are explicitly listed
in \File{equations/polynomial.json}. Neither file contains a factorial,
variable exponent, binomial operator, bounded quantifier, unspecified
Diophantine representation, external constant port, or oracle call.
The files are two exact representations of the same finite polynomial.

\begin{theorem}[Fixed consistency equation]\label{thm:main}
For the polynomial specified by \eqref{eq:final} and the accompanying
coefficient list, ZFC proves
\[
 (\exists\mathbf z\in\ZZ^{1873}\;P(\mathbf z)=0)
 \quad\longleftrightarrow\quad \Refut_{\mathrm{ZFC}}.
\]
Here $\Refut_{\mathrm{ZFC}}$ means existence of a finite syntactic derivation
of falsity in the repository's stated ZFC proof system. Consequently ZFC
proves the no-zero/$\Con(\mathrm{ZFC})$ biconditional.
\end{theorem}

The proof composes the inherited C222 calibration with the new arithmetic
transformations proved below. It does not assume consistency, the existence
of a standard model of ZFC, or a bound on a possible refutation.
The output is a mathematical proof with exact finite computational
certificates. The complete composed theorem has not been checked by a
proof assistant during this work.

\subsection{Two supplied quartics}

\begin{center}\small
\begin{tabular}{lrr}
\toprule
Measure & Primary & Smaller coefficient/term variant\\
\midrule
Integer variables & 1,873 & 3,292\\
Degree & 4 & 4\\
Quadratic residuals & 1,468 & 2,887\\
Collected monomials & 77,523 & 11,218\\
Coefficient height & $1{,}323{,}565{,}948{,}780{,}873{,}728$ & $14{,}869{,}246$\\
Coefficient bit length & 61 & 24\\
\bottomrule
\end{tabular}
\end{center}

The second equation is in \File{equations_small_height/}. It keeps more
affine arithmetic registers, so its individual residuals have at most
five monomials. Both variants have the same semantic certificate and
the same consistency theorem. Their differences are exact choices of
intermediate registers, with unique integral restoration.

\subsection{Comparison with the inspected repository}

The comparison baseline is the actual repository at the commit printed
on the title page, including the then-new
\File{ConZFC_Direct_Projection_7617.zip}, rather than only the older
manuscript-collection README~\cite{repo,projection,latest}.

\begin{center}\small
\begin{tabular}{lrrr}
\toprule
Direct integer construction & Variables & Degree & Monomials\\
\midrule
Earlier direct C222 round 12 & 27,916 & 4 & 758,275\\
Shared CRT, prime-power lift & 8,619 & 4 & 39,811\\
Incoming eight-column projection & 7,617 & 4 & 148,220\\
This work, primary & \textbf{1,873} & 4 & 77,523\\
This work, retained affine registers & 3,292 & 4 & \textbf{11,218}\\
\bottomrule
\end{tabular}
\end{center}

The primary count is about $75.4\%$ below 7,617 and $78.3\%$ below 8,619.
The latter parent has fewer monomials than the primary; the second supplied
variant addresses that tradeoff. These comparisons are not uniform claims
of improvement in every measure.

The repository also records nine-positive-variable and ten-integer-variable
high-degree equations, and quartics with 37 integer variables using the
native coefficient construction~\cite{size}. Thus 1,873 is not a record
variable count among all available routes. The contribution here is a
substantially smaller direct proof-calculus equation with fully expanded,
moderate-size literal coefficients and a reusable simultaneous arithmetic
theorem. Historical priority for the underlying idea is not asserted.

\section{Pell notation and the inherited finite theorems}

For $a\ge2$ define
\begin{align*}
 \chi_a(0)&=1,&\psi_a(0)&=0,\\
 \chi_a(j+1)&=a\chi_a(j)+(a^2-1)\psi_a(j),&
 \psi_a(j+1)&=\chi_a(j)+a\psi_a(j).
\end{align*}
Equivalently,
$(a+\sqrt{a^2-1})^j=\chi_a(j)+\psi_a(j)\sqrt{a^2-1}$.
All quantities used in the proof can instead be defined by the displayed
integer recurrence. In particular
\[
 \chi_a(j)^2-(a^2-1)\psi_a(j)^2=1.
\]
For $j>0$, both coordinates are positive, $\psi_a(j)\ge j$, and
$\chi_a(j)\ge a^j$. Every positive-abscissa nonnegative-ordinate solution
of this Pell equation is one of these pairs.

The inherited arithmetic input is explicit: the actual statements and
proofs of \File{Pell.matiyasevic} and \File{Pell.eq_pow_of_pell} in
Mathlib~\cite{pell}, together with the repository's literal polynomial
interfaces \File{PellRepresentation.lean} and
\File{PowerRepresentation.lean}. Copies are included in
\File{provenance/}. The inspected Mathlib blob is
\texttt{db9c757dd7cac29260397150fd190d72715ea3f8}.

\begin{proposition}[Positive power criterion]\label{prop:power}
For $n,k>0$ and $m\ge0$, $m=n^k$ is equivalent to the existence of natural
$w,a,t,z$ with $a\ge2$ satisfying
\begin{gather*}
 \chi_a(k)\equiv\psi_a(k)(a-n)+m\pmod t,\qquad
 2an=t+n^2+1,\\
 m<t,\qquad n\le w,\quad k\le w,\qquad
 a^2-((w+1)^2-1)(wz)^2=1.
\end{gather*}
In the formal source $a-n$ is natural subtraction. Our construction proves
$a>n$ before applying this proposition.
\end{proposition}

The usual finite index certificate for
$(x,y)=(\chi_a(k),\psi_a(k))$ uses three Pell norms and four congruences.
The next section proves the signed variant needed here, with all domain
conditions stated explicitly. This removes a genuine arithmetic
restriction from the compiler; it is not just a relabeling of variables.

\section{Signed recovery of a positive Pell index}\label{sec:signed}

We use two precise properties of Pell sequences~\cite{pell}:
\begin{enumerate}
\item If $\psi_a(i)^2\mid\psi_a(n)$, then $\psi_a(i)\mid n$.
This is \File{dvd_of_ysq_dvd} in the supplied source.
\item If $0<i\le n$ and
$\chi_a(h)\equiv\chi_a(i)\pmod{\chi_a(n)}$, then
$h\equiv i$ or $h\equiv-i\pmod{4n}$.
This is \File{modEq_of_xn_modEq}; its positivity hypothesis excludes
the exceptional zero-index coincidence in the preceding auxiliary lemmas.
\end{enumerate}
The addition formula also gives
\begin{equation}\label{eq:signshift}
 \chi_a(h+2n)\equiv-\chi_a(h)\pmod{\chi_a(n)}.
\end{equation}
Indeed $\chi_a(2n)=2\chi_a(n)^2-1$ and
$\psi_a(2n)=2\chi_a(n)\psi_a(n)$. These facts specify exactly what is
needed from the modular Pell theory; no empirical pattern is substituted
for an index theorem.

\begin{lemma}[Signed index recovery]\label{lem:signed}
Let $A,B,U,X,Y,k\in\NN$, with $A,B\ge2$, $U\ge2$, $X>0$ and $1\le k\le Y$.
Let $V,S,R\in\ZZ$. Assume
\begin{gather}
 X^2-(A^2-1)Y^2=1,\quad U^2-(A^2-1)V^2=1,
 \quad S^2-(B^2-1)R^2=1,\label{eq:three-norms}\\
 Y^2\mid V,\quad B\equiv A\pmod U,\quad B\equiv1\pmod{4Y},
 \quad S\equiv X\pmod U,\quad R\equiv k\pmod{4Y}.\label{eq:signed-cong}
\end{gather}
Then $X=\chi_A(k)$ and $Y=\psi_A(k)$. No sign hypothesis on $S$ or $R$
is needed.
\end{lemma}

\begin{proof}
Since $U\ge2$, the second norm forces $V\ne0$. Pell classification gives
indices $i,n,h$ with
\[
 (X,Y)=(\chi_A(i),\psi_A(i)),\quad
 (U,|V|)=(\chi_A(n),\psi_A(n)),\quad
 (|S|,|R|)=(\chi_B(h),\psi_B(h)).
\]
We have $i,n>0$. From $Y^2\mid|V|$ we obtain $Y\le|V|$, hence $i\le n$
by strict monotonicity; also $Y\mid n$ by property 1 above.

The coordinate polynomials preserve congruences in their base. Consequently
$B\equiv A\pmod U$ and the congruence for $S$ imply
\[
 \chi_A(h)\equiv \pm\chi_A(i)\pmod U.
\]
In the positive-sign case property 2 gives $h\equiv\pm i\pmod{4n}$.
In the negative-sign case apply \eqref{eq:signshift} and property 2 to
$h+2n$. In both cases
\begin{equation}\label{eq:h-i}
 h\equiv\pm i\pmod{2n}.
\end{equation}
This step is why taking absolute values alone would have been insufficient:
the possible change of sign is absorbed in the index modulus.

At base 1 the recurrence has ordinate $h$. Thus
$B\equiv1\pmod{4Y}$ implies $\psi_B(h)\equiv h\pmod{4Y}$.
The signed congruence for $R$ therefore gives
\[
 h\equiv\pm k\pmod{4Y}.
\]
Reduce this congruence and \eqref{eq:h-i} modulo $2Y$, which divides
$2n$. It follows that $i\equiv\pm k\pmod{2Y}$.
Finally $1\le i\le Y$ because $\psi_A(i)\ge i$, and $1\le k\le Y$
by hypothesis. The difference congruence gives $i=k$; the sum congruence
can only give $i+k=2Y$, which also forces $i=k=Y$.
This includes the boundary cases $Y=1$ and $i=n$.
\end{proof}

In the final polynomial, this lemma permits the two variables in
$S=X+Us$ and $R=k+4Yr$ to range over all integers. For sixteen seed
indices it removes 32 nonnegativity equations and their 96 auxiliary
integer coordinates. All remaining necessary sign restrictions are kept.

\section{One simultaneous certificate for all powers}\label{sec:batch}

Fix a nonempty finite list of positive input/output triples
$(n_i,k_i,m_i)$, $1\le i\le q$. Initially group identical exponents into
$r$ classes, with exponent $e_j$ and class map $j(i)$.
Introduce natural variables $g,u,b,x_j,y_j$ and arbitrary integer
variables $z,v,\lambda_B,s_j,r_j,\lambda_i$. Define polynomial
abbreviations
\begin{align}
 W&=\sum_{i=1}^q(n_i+k_i),&
 A&=2+\sum_{i=1}^q m_i+g,&U&=u+2,\label{eq:shared-abbr1}\\
 X_j&=x_j+1,&Y_j&=e_j+y_j,&
 P_*&=\prod_{j=1}^rY_j,\label{eq:shared-abbr2}\\
 V&=4P_*^2v,&B&=1+4P_*(b+1),&
 T_i&=2An_i-n_i^2-1,\label{eq:shared-abbr3}\\
 S_j&=X_j+Us_j,&R_j&=e_j+4Y_jr_j.\label{eq:shared-abbr4}
\end{align}
The star distinguishes $P_*$ from the final quartic $P$.
These abbreviations introduce no quantified variables.

Impose
\begin{align}
 A^2-((W+1)^2-1)(Wz)^2&=1,\label{eq:global1}\\
 U^2-(A^2-1)V^2&=1,\label{eq:global2}\\
 B-A&=U\lambda_B,\label{eq:global3}\\
 X_j^2-(A^2-1)Y_j^2&=1 &&(1\le j\le r),\label{eq:local1}\\
 S_j^2-(B^2-1)R_j^2&=1 &&(1\le j\le r),\label{eq:local2}\\
 X_{j(i)}+Y_{j(i)}n_i-Y_{j(i)}A-m_i&=T_i\lambda_i
 &&(1\le i\le q).\label{eq:power-row}
\end{align}
An emitted congruence may use the negative of the displayed quotient;
since these quotients are unrestricted integers this is an exact
coordinate change.

\begin{theorem}[Simultaneous power graph]\label{thm:batch}
For specified positive triples $(n_i,k_i,m_i)$, all equalities
$m_i=n_i^{k_i}$ hold if and only if
\eqref{eq:global1}--\eqref{eq:power-row} have auxiliary values in the
stated domains.
\end{theorem}

\subsection{Soundness and the single common output bound}

First $W\ge2$, $A\ge2$, $U\ge2$, $Y_j\ge e_j\ge1$,
$P_*\ge1$, and $B\ge5$. If $z=0$, \eqref{eq:global1} would give
$A^2=1$. Therefore $z^2\ge1$, and rearrangement gives
\[
 A^2=1+W^3(W+2)z^2\ge1+W(W+2)=(W+1)^2.
\]
Thus $A\ge W+1>n_i$ for every call. The exact identity
\begin{equation}\label{eq:output-bound}
 T_i-A=(n_i-1)(2A-n_i-1)+(A-2)
\end{equation}
shows $T_i\ge A>m_i$. This argument uses only the displayed polynomial
system and the required domains. In particular it does not first assume
that the candidate power outputs are correct. One global gap $g$ has
replaced all $q$ private strict output-bound slacks.

Equation \eqref{eq:global2} forces $v\ne0$. Since each $Y_j$ divides
$P_*$, we have $Y_j^2\mid V$ and $B\equiv1\pmod{4Y_j}$.
The definitions of $S_j,R_j$ and \eqref{eq:global3} give the remaining
congruences in Lemma~\ref{lem:signed}. That lemma identifies every
$(X_j,Y_j)$ with the true pair at $e_j$ in base $A$.

Now \eqref{eq:power-row}, \eqref{eq:output-bound}, $n_i,k_i\le W$,
and \eqref{eq:global1} with $|z|$ satisfy
Proposition~\ref{prop:power}. The natural subtraction there agrees
with integer subtraction because $A>n_i$. Hence every $m_i=n_i^{k_i}$.

\subsection{Completeness without incompatible shared choices}

Suppose all candidate outputs are correct. Fix all input/output values
first. This is legitimate even if the surrounding arithmetic system
relates different calls; a simultaneous true tuple is already given.

\paragraph{Stage 1: the common bootstrap.}
At base $W+1$ the recurrence gives
$\psi_{W+1}(\ell)\equiv\ell\pmod W$.
Choose a sufficiently large positive multiple $\ell=Wh$ and set
\[
 A=\chi_{W+1}(\ell)\ge2+\sum_i m_i,\qquad
 z=\psi_{W+1}(\ell)/W,\qquad
 g=A-2-\sum_i m_i.
\]
The sequence is increasing and unbounded, so such an $h$ exists.
All choices are finite natural-number constructions, and
\eqref{eq:global1} holds. Define
$X_j=\chi_A(e_j)$ and $Y_j=\psi_A(e_j)$.
Then $x_j=X_j-1$ and $y_j=Y_j-e_j$ are nonnegative.

\paragraph{Stage 2: the common auxiliary Pell solution.}
Form $P_*\ge1$ and $K=4P_*^2$. The matrix
\[
 M_A=\begin{pmatrix} A&A^2-1\\1&A\end{pmatrix}
\]
has determinant 1, so it is invertible modulo $K$. In the finite group
of invertible $2\times2$ matrices modulo $K$, some positive power is
the identity. For such an exponent $d$ put
\[
 U=\chi_A(d),\qquad V=\psi_A(d).
\]
Then $U\equiv1\pmod K$, $K\mid V$, $V>0$, and $U\ge A\ge2$.
Set $u=U-2$ and $v=V/K$. In particular
$\gcd(U,4P_*)=1$.

\paragraph{Stage 3: the common CRT base.}
Choose, by CRT and addition of a positive period, a number
\[
 B>A,\qquad B\equiv A\pmod U,\qquad B\equiv1\pmod{4P_*}.
\]
Take it large enough that $b=(B-1)/(4P_*)-1\ge0$.
For each exponent put $S_j=\chi_B(e_j)$ and $R_j=\psi_B(e_j)$.
The recurrence modulo $U$ gives $S_j\equiv X_j\pmod U$;
the recurrence at base 1 modulo $4Y_j$ gives
$R_j\equiv e_j\pmod{4Y_j}$. Thus the required integer $s_j,r_j$
exist. The local and global norms hold by construction.

Finally, the Pell power congruence gives
\[
 X_j+Y_jn_i-Y_jA-n_i^{k_i}\equiv0\pmod{T_i}.
\]
It follows either directly by expanding powers in the quadratic ring,
or from \File{x_sub_y_dvd_pow} in the supplied arithmetic source.
The positive integer $T_i$ therefore supplies $\lambda_i$.
This proves completeness and Theorem~\ref{thm:batch}.

\begin{remark}[Why the argument can share all calls]
There is no circular witness bound. The common bootstrap is chosen after
the true input/output tuple. The product of ordinates is chosen after
the bootstrap. The auxiliary Pell solution is chosen after that product,
and the CRT base after the auxiliary solution. None of these private
choices changes a previously fixed input/output value.
\end{remark}

\section{Sharing neighboring exponents}\label{sec:recurrence}

The actual arithmetic source has forty distinct power calls but only
nineteen distinct exponent expressions. Three are immediate shifts of
two others:
\begin{center}
\begin{tabular}{ccc}
\toprule
Derived exponent & Seed exponent & Shift\\
\midrule
$L+2$ & $L$ & 2\\
$N+2$ & $N$ & 2\\
$N+1$ & $N$ & 1\\
\bottomrule
\end{tabular}
\end{center}
Here $L$ is the actual factorial scale argument and $N$ the positive
proof-row count. The emitted builder proves these identities by exact
affine normalization over nonlinear arithmetic nodes. It does not rely
on fixed node numbers or an assumed semantic equality.

Let $(X_0,Y_0)=(\chi_A(k),\psi_A(k))$ and $D=A^2-1$. Compute
\begin{align*}
 X_1&=AX_0+DY_0,&Y_1&=X_0+AY_0,\\
 X_2&=2AX_1-X_0,&Y_2&=2AY_1-Y_0.
\end{align*}
These are exactly the pairs at $k+1$ and $k+2$.
The second-order expressions use fewer multiplications than a second
independent first-order step. They also equal the nonnegative formulas
\begin{align*}
 X_2&=(2A^2-1)X_0+2A(A^2-1)Y_0,\\
 Y_2&=2AX_0+(2A^2-1)Y_0.
\end{align*}
Thus the signed intermediate registers do not create additional
semantic cases.

Only sixteen seed pairs need Lemma~\ref{lem:signed}. Define $P_*$ as
the product of just their ordinates. Once the seed pairs are identified,
the recurrences identify all derived pairs. The power criterion then
applies to every call. No divisibility or index certificate is needed
for a pair already defined by an exact recurrence.

Before applying the signed index refinement, this removes twelve
semantic inputs and 36 domain witnesses, plus a net 25 nonlinear
registers: 73 integer variables, taking 2,042 to 1,969. The signed
refinement then removes 96 more coordinates, taking 1,969 to 1,873.
These counts state the order of comparison; the recurrence theorem
itself applies equally to the signed certificate.

\section{The fixed logical input: a direct numerical calculus}

The source is the fixed C222 code calculus from ProveIt~\cite{calculus}.
The new equation does not enumerate a Turing-machine execution.
The relevant inherited interface is the machine-independent theorem
\File{KEndpoint.c222_iff_zf_refutation}, composed with
\File{ConstructibleChoice.zf_refutation_iff_zfc_refutation}.
Both theorem statements and their source provenance accompany this work.
The constructible-model argument is the established logical calibration,
not a new consequence of the arithmetic batching lemma.

For clarity the numerical constructors can be printed compactly. Put
\begin{align*}
 \I(a,b)&=2(a+b)(a+b+1)+4b,&
 \E(a,b)&=\I(a,b)+1,\\
 \Q(a,b)&=\I(a,b)+2,&
 \J(a,b)&=\I(a,b)+3,\\
 \Neg(a)&=\I(a,1),&
 \Ex(a,b)&=\Neg(\Q(a,\Neg(b))).
\end{align*}
Code 0 is truth and code 1 is falsity. The eight ordinary axiom
conclusions are
\begin{align*}
 \rho_1(a,b)&=\Ex(a,\J(a,b)),\\
 \rho_2(a,b,c)&=\I(\Q(a,\I(b,c)),\I(\Q(a,b),\Q(a,c))),\\
 \rho_3(a,b,c)&=\I(\J(a,b),\I(\E(c,a),\E(c,b))),\\
 \rho_4(a,b,c)&=\I(\Q(a,\Q(b,c)),\Q(b,\Q(a,c))),\\
 \rho_5(a,b,c)&=\I(\J(c,a),\I(\E(c,b),\E(a,b))),\\
 \rho_6(a,b)&=\I(\Ex(a,\Q(a,b)),b),\\
 G(a,b,c)&=\I(\J(c,a),\I(\J(c,b),\Q(\E(a,b),\J(a,b)))),\\
 \rho_8(a,b,c,d)&=\I(\I(\I(c,d),a),\I(\I(a,c),\I(b,c))).
\end{align*}
The remaining axiom conclusion is the seed bundle. With
\[
 \mathrm{req}(b,k,m)=\Q(b,\I(m,\Neg(\Q(b,\I(m,\Neg(\E(b,k))))))).
\]
To keep the closure constructor readable, introduce three exact abbreviations:
\begin{align*}
 \mathrm{sb}(a,b,k)&=\Q(k,\I(\E(k,a),\E(k,b))),\\
 \mathrm{c}_1(a,b,k,m)&=\Q(a,\I(\mathrm{sb}(a,b,k),\mathrm{req}(b,k,m))),\\
 \mathrm{c}_2(a,b,k,m)&=\Q(a,\I(m,\Neg(\mathrm{c}_1(a,b,k,m)))),\\
 \mathrm{cl}(a,b,k,m)&=\Q(b,\I(\E(b,k),\Neg(\mathrm{c}_2(a,b,k,m)))).
\end{align*}
define
\begin{align*}
 \mathrm{nks}(a,b,c)&=\Q(c,\I(\mathrm{cl}(a,b,c,c),\Neg(\mathrm{req}(b,c,c)))),\\
 \mathrm{vr}(a,b,c)&=\I(\E(a,b),\Q(c,\E(a,b))),\\
 \mathrm{bundle}(a,b,c)&=\Neg(\I(\mathrm{vr}(a,b,c),\mathrm{nks}(a,b,c))).
\end{align*}
Its guard is $c\ne a,b$. The displayed constructors and the literal
JSON constructor graph have been independently checked against the
pinned code definitions.

There are twelve single-conclusion rule branches: truth; these nine
axiom conclusions; modus ponens with premises $a,\I(a,b)$ and conclusion
$b$; and generalization with premise $b$ and conclusion $\Q(a,b)$.
Unused premises are code 0. The bundle guard has the exact existential
natural form
\[
 ((c-a)(c-b))^2=(d+1)^2.
\]
It holds for some $d\in\NN$ exactly when $c\ne a$ and $c\ne b$.
All conclusion constructors are polynomials with nonnegative coefficients.

The calibration we inherit states
\begin{equation}\label{eq:calibration}
 \mathrm{C222}(1)\quad\longleftrightarrow\quad\Refut_{\mathrm{ZFC}}.
\end{equation}
This is the only logical-strength interface needed by the new direct
compiler. The arithmetic proof does not require a semantic interpretation
for every possible intermediate natural code.

\section{From a finite derivation to eight bounded columns}

We summarize the inherited eight-column projection in sufficient detail
to identify every proof obligation~\cite{projection}. The complete source
proof is retained in \File{proofs/SOURCE_PROJECTION.md}.

\subsection{Padding and a literal label majorant}
For each of the twelve branches let $p_b,q_b,v_b$ be its two premise
labels and conclusion after substituting its four natural parameters
$a,b,c,d$. Choose a proof-tree expansion deriving code 1. Pad its list of
rule instances with truth steps until its length $N\ge1$ bounds all four
parameters of every step. Define the explicit nonnegative polynomial
\[
 H(N)=\sum_b\bigl(p_b(N,N,N,N)+q_b(N,N,N,N)+v_b(N,N,N,N)\bigr).
\]
It bounds every label in the list. Put
\[
 M_0=2^{2H(N)N+1},\qquad M=M_0+1.
\]
This power is one of the forty calls later replaced by the simultaneous
polynomial certificate.

Use an accumulator with $S_0=M+1$, $S_N=1$, and updates
\begin{equation}\label{eq:accumulator}
 S_{i+1}(1+v_iM)=S_i(1+p_iM)(1+q_iM).
\end{equation}
For a proof tree, $S_i$ is the product of the factors for outstanding
nonzero obligations. Expanding one obligation gives the displayed update;
the empty product at the end is 1.

Conversely, all factors are positive and $S_0>0$, so induction gives
$S_i>0$ for every $i$. Telescoping and cancellation in the updates
and endpoints therefore imply
\[
 \prod_i(1+v_iM)=(1+M)\prod_i(1+p_iM)(1+q_iM).
\]
The coefficient magnitudes of the difference of these two polynomials
in a formal argument are at most $M_0<M$. This follows by evaluating
the nonnegative coefficient sums at 1 and using $H+1\le2^H$.
A least nonzero coefficient cannot vanish on evaluation at $M$:
reduction modulo $M$ would make it divisible by $M$. Thus the two
polynomials are identical.

Comparing multiplicities at the distinct roots $-1/j$ for positive labels
$j$ gives the Horn balance equation
\[
 \sum_i(e_{v_i}-e_{p_i}-e_{q_i})=e_1,\qquad e_0=0.
\]
Let $D$ be the least closure of the selected rules with truth available.
If 1 were outside $D$, the sum of balance coordinates outside $D$ would
be nonpositive on each rule: an outside conclusion must have an outside
premise. Its value on $e_1$ is 1, a contradiction. Therefore code 1 is
derivable. This proof allows cycles, repeated premises and self-loops.

\subsection{The two adjacent beta digits}
Encode the finite accumulator by the ordinary remainder relation
\[
 x=q_i(1+(i+1)y)+S_i.
\]
Choose $y$ above every $S_i$ and divisible by $N!$; the moduli for
$0\le i\le N$ are pairwise coprime, so CRT supplies $x$.
Let $U_0$ bound all the finitely many quotients.
At row $i$ the only bounded local coordinates are
\[
 (a,b,c,d,S_i,q_i,S_{i+1},q_{i+1})
 \quad\text{with bounds}\quad (N,N,N,N,y,U_0,y,U_0).
\]
The bound $S_i\le y<1+(i+1)y$ already proves the strict remainder
condition. No separate local remainder-bound slack is needed.

For each branch form the square of its accumulator residual and, for the
bundle branch, add the square of its guard residual. Multiply all twelve
branch sums, then add the squares of the two beta reconstruction
residuals. Call the resulting fixed integer polynomial $R(i,\mathbf w)$.
Its zero set is exactly the row relation. The truth, modus ponens and
generalization rows use the literal premise/conclusion substitutions
above. No independent constructor coordinates cross the bounded quantifier.

\section{Signed row evaluation and the absolute majorant}

Represent $R$ by a fixed arithmetic graph with operations $+,-,\times$.
Its absolute majorant $R^+$ replaces every subtraction by addition and
every literal constant by its absolute value. It is a polynomial graph
with nonnegative coefficients.

\begin{lemma}[Circuit majorant]
If $0\le i\le N$ and $0\le w_j\le u_j$, while the other natural
parameters are held fixed, then
\[
 |R(i,\mathbf w)|\le R^+(N,\mathbf u).
\]
\end{lemma}
\begin{proof}
Induct through the graph, using the triangle inequality at addition and
subtraction, multiplicativity of absolute value at multiplication, and
monotonicity of nonnegative arithmetic at the input bounds.
\end{proof}

We evaluate the signed graph directly for the row congruence and use
only its absolute majorant for the range argument. Expanding both halves
of a positive/negative pair representation at both interfaces was
unnecessary. This saves 80 final integer coordinates compared with the
earlier paired evaluation in the synchronized construction. The finite
source checker proves the complete signed-row identity compositionally.

\section{A bounded-family certificate with lifted CRT periods}\label{sec:crt}

Fix the outer parameters and the bounds $\mathbf u$. Set
\[
 L=R^+(N,\mathbf u)+N,\quad T=L!,\quad
 d_i=1+(i+1)T,\quad C=\prod_{i<N}d_i.
\]
The source has $N\ge1$, hence $L\ge1$. The following finite arithmetic
conditions express the whole bounded family:
\begin{equation}\label{eq:range}
 C\mid\binom{A_j}{u_j+1}\;(1\le j\le8),\qquad
 C\mid1+(c+1)T,\qquad R(c,\mathbf A)\equiv0\pmod C.
\end{equation}
The encoded columns are parameterized as $A_j=\alpha_j+1$, allowing
zero decoded values while making the arithmetic upper arguments positive.

\begin{proposition}[Bounded row equivalence]
Conditions \eqref{eq:range} have natural encoded columns and index if
and only if every $i<N$ has a natural row witness bounded by
$\mathbf u$ with $R(i,\mathbf w)=0$.
\end{proposition}

\begin{proof}[Soundness]
Choose a prime $p\mid d_i$. Since $\gcd(d_i,T)=1$, it does not divide
$L!$, so $p>L$. The index condition gives $c\equiv i\pmod p$.
Multiplying a range divisibility by $(u_j+1)!$ and using the integer
falling-factorial identity gives
\[
 p\mid\prod_{t=0}^{u_j}(A_j-t).
\]
Some $w_j\in[0,u_j]$ therefore satisfies $A_j\equiv w_j\pmod p$.
The row congruence yields $R(i,\mathbf w)\equiv0\pmod p$.
The majorant bounds its absolute value by $L<p$, so it is zero.
There is no cancellation of a possibly noninvertible factorial here.
\end{proof}

\begin{proof}[Completeness]
The $d_i$ are pairwise coprime: a common prime would divide a nonzero
difference of indices smaller than $N$ and hence would divide $L!$,
which contradicts $d_i\equiv1\pmod{L!}$.
For each column put $F_j=(u_j+1)!$ and lift the periods to
\[
 D_{ij}=\prod_{p\mid d_i}p^{v_p(d_i)+v_p(F_j)}.
\]
For fixed $j$ these moduli are still pairwise coprime. Choose
$A_j\equiv w_{ij}\pmod{D_{ij}}$ for all $i$, and shift by their product
to make $A_j>u_j$. The numerator of the falling factorial has a factor
with enough valuation to absorb both $d_i$ and $F_j$. Division by
$F_j$ therefore leaves $d_i\mid\binom{A_j}{u_j+1}$.
Choose the encoded index independently modulo the original $d_i$.
The strengthened column congruences imply the original row congruences.
All of \eqref{eq:range} follows.
\end{proof}

The lifted periods occur only in the proof of existence; they add no
variable or prime-factorization operator to the emitted polynomial.
A shift by the unlifted product is insufficient: for example
$15\mid\binom0{5}$ but $15\nmid\binom{15}{5}=3003$.
This is why merely requesting positive encoded columns without a
compatible-period proof would be invalid.

\section{Eliminating the remaining arithmetic operations}

The CRT product has the explicit formula
\begin{align*}
 Q&=(1+(N+1)T)^N+2,\\
 C&=\bigl(T^N N!\binom{Q+N}{N}\bigr)\bmod(TQ-1).
\end{align*}
Modulo $TQ-1$, the numerator is congruent to
$\prod_{j=1}^N(1+jT)$ because $TQ\equiv1$.
Also that product is at most $(1+NT)^N<Q-1\le TQ-1$,
so the remainder is exactly the desired positive product.

For a binomial call with positive $n,k$, let $q=2^n+1$,
$D=q^k$ and $F=(q+1)^n$. Use
\[
 F=D(qh+b)+r,\qquad D=r+r_s+1,\qquad q=b+b_s+1,
\]
with $b,r,r_s,b_s\in\NN$ and $h\in\ZZ$.
If $h<0$, then $qh+b<0$ and the first equation would give $F<0$.
Thus $h\ge0$ is forced. The displayed relations say that $b$ is the
base-$q$ digit of $\lfloor F/D\rfloor$.
Every binomial coefficient of order $n$ is smaller than $q$, so
$b=\binom nk$, including the case $k>n$.

For a factorial argument $n\ge1$, put
\[
 R_0=(n+2)^{n+2},\qquad c=\binom{R_0+n}{n}.
\]
The explicit factorial identity in the supplied source is
\[
 n!=\left\lfloor\frac{(R_0+1)^n+c-1}{c}\right\rfloor.
\]
The quadratic-graph interface imposes
\[
 (R_0+1)^n+c=cf+r+1,\qquad c=r+s+1,
\]
with natural remainder and slack. We parameterize $f=f_0+1$.
This is valid because every factorial is positive.

These constructions leave forty power calls, eleven binomial calls,
two factorial calls and three ordinary division calls. Each power
output is parameterized as $m=m_0+1$. The actual graph then establishes
\emph{all forty base and exponent positivity conditions syntactically},
before assuming any power or factorial equality. This avoids a circular
application of the simultaneous theorem.

The nine coordinate/index divisibility quotients are unrestricted:
their natural numerators and the now-proved positive CRT product force
them to be nonnegative. Similarly the three ordinary Euclidean quotients
are forced nonnegative by a natural dividend and a bounded natural
remainder. These facts justify the omitted guards. The checkers verify
that these inferred signs do not feed a power-argument positivity proof.

\section{From the mixed arithmetic graph to a quartic over integers}

For every semantic input $x$ still required to be nonnegative, add three
fresh unrestricted integers $a,b,c$ and the quadratic residual
\begin{equation}\label{eq:nonneg}
 x-a^2-b^2-c(c+1)=0.
\end{equation}
The right side of $x=a^2+b^2+c(c+1)$ is always nonnegative.
Conversely, Legendre's three-square theorem applies to $4x+1$ for every
$x\ge0$: it is not of the form $4^a(8b+7)$.
Modulo four, precisely one of the three representing squares is odd.
Writing the even squares as $(2a)^2,(2b)^2$ and the odd one as
$(2c+1)^2$ yields \eqref{eq:nonneg}.

All other arithmetic registers range freely over $\ZZ$. Substitute
addition, subtraction and scalar multiplication as affine forms.
For every genuine product introduce a fresh variable $z$ and the
quadratic equation $z-L_1L_2=0$, where $L_1,L_2$ are the current affine
forms. If an affine form exceeds sixteen monomials, retain it as a
register with a linear defining equation. The smaller-height variant
uses a threshold of one.

Given semantic coordinates, topological evaluation supplies unique
integer values for every retained register. Conversely, the defining
equations force those values in order. Hence this flattening preserves
the complete integer zero set up to unique register restoration.
The nonnegativity representations are existentially surjective, as
proved above. Finally a sum of integer squares vanishes exactly when
all its summands vanish.

\begin{center}\small
\begin{tabular}{lrr}
\toprule
Component & Variables & Residuals\\
\midrule
Semantic inputs & 245 & --\\
Source arithmetic equalities & -- & 57\\
Shared power-certificate equalities & -- & 75\\
Nonlinear defining registers & 1,153 & 1,153\\
Retained affine registers & 37 & 37\\
146 nonnegativity representations & 438 & 146\\
\midrule
Total & \textbf{1,873} & \textbf{1,468}\\
\bottomrule
\end{tabular}
\end{center}
Of the 245 semantic inputs, 146 require nonnegativity and 99 are
unrestricted. The latter include the 32 newly unguarded Pell progression
coordinates. The degree is at most four by construction, and exactly
four because a fresh guard coordinate has fourth-power coefficient 1
and occurs in no other residual. The collected expansion independently
confirms the degree, every coefficient and all-variable liveness.

\section{Proof of the consistency theorem}

Every step in the following chain has been given both existence directions:
\begin{align*}
 P(\mathbf z)=0\text{ for some }\mathbf z\in\ZZ^{1873}
 &\iff \text{the mixed arithmetic certificate has a solution}\\
 &\iff \text{the bounded family of direct proof rows holds}\\
 &\iff \mathrm{C222}(1)\\
 &\iff \Refut_{\mathrm{ZFC}}.
\end{align*}
The first equivalence uses the register restoration and
\eqref{eq:nonneg}. The second uses the signed shared Pell theorem,
recurrence substitution, exact binomial/factorial/division graphs and
the lifted CRT theorem. The third is the padded Horn-flow construction.
The last is the pinned logical calibration \eqref{eq:calibration}.

All new arguments are finite algebra, elementary integer arithmetic,
Pell recurrence theory, finite-group reasoning, CRT and finite proof
combinatorics. They can be carried out inside ZFC over its own internal
$\omega$. Neither direction is a meta-level appeal to the truth of
consistency. The chain proves Theorem~\ref{thm:main}; classical negation
gives its no-zero/$\Con(\mathrm{ZFC})$ consequence.

The attached Lean files document exact inherited theorem interfaces.
They are not advertised as a standalone compilable Lean project, and
we have not emitted an internal proof object for the new composed
biconditional. This distinction does not turn the ordinary argument
into a conjecture; it states the form of the evidence being delivered.

\section{Reproduction, exact checks and provenance}

All Python programs use only the standard library. From the package root:
\begin{verbatim}
python3 build.py
python3 build.py --affine-cap 1 --output equations_small_height
python3 verify_release.py
\end{verbatim}
The release verifier reconstructs the two equations in temporary
directories, compares every generated byte, and runs the independent
checks in normal and optimized Python. The mathematical proofs remain
necessary: evaluating finitely many integer tuples could not establish
the infinite zero-set equivalence.

The independent checks have complementary scopes:
\begin{itemize}
\item \File{verify_source.py}: every one of the twelve rule branches,
the paired and signed local row, the absolute majorant, the full
heterogeneous bounds, both beta endpoints, all 57 source equalities,
and every arithmetic-call interface. It checks unconditional positivity
of all forty power bases and exponents. Its companion
\File{verify_source_negative.py} rejects three corruptions of the
bounds, majorant and horizon domain.
\item \File{verify_batch.py}: all common definitions, seed and derived
Pell coordinates, every common/local norm, all forty power congruences,
the exact domain choices and the entire 75-row batch. Its negative
controls corrupt meaningful formulas and required domain guards.
\item \File{verify_recurrence.py}: the free polynomial recurrence
identities, the three actual exponent differences and the exact product
of sixteen seed ordinates.
\item \File{verify_polynomial.py}: complete pruning, every affine and
nonlinear register equation, all nonnegativity rows, and exact
collection of every coefficient in the sum of squares. A coefficient
change and a guard-square change are detected.
\end{itemize}

The checkers do not import the equation builder. The coefficient checker
uses symmetric convolution, independently of the emitter's ordered
double multiplication loop. Checks use explicit exceptions rather than
assertions that disappear under \texttt{python3 -O}.

\subsection{The finite encoding of the equation}
The polynomial JSON uses a list of pairs \File{[coefficient,monomial]}.
A monomial is a sorted list of zero-based variable indices, with repeats
representing powers. For example \File{[7,[0,0,3]]} means $7z_0^2z_3$,
and the empty list denotes 1. Coefficients are decimal integer literals.
The residual JSON uses the same convention. These rules specify the
entire equation without requiring the reader to execute the generator.

The primary polynomial's SHA-256 is
\begin{quote}\small\ttfamily
18d09e54a3d66ed33422560a35724a5073\\
fd1ca1a2a76bfaf16e634cb019ea1d
\end{quote}
The line break is typographical; concatenate the two lines.
The release checksum manifest covers the immutable source, equation,
proof and document artifacts. The regenerable \File{checks/} receipts
and the manifest itself are excluded.

\subsection{Attribution and proof boundary}
The eight-column source, Horn-flow argument, C222 calibration, prime-power
CRT lift, merged binomial/factorial formulas and three-square integer
conversion have explicit predecessors in the inspected repository.
The new contributions in this package are their concrete synthesis with
the shared bootstrap/CRT power certificate, the one-gap output bound,
the signed positive-index recovery theorem, exact neighboring-exponent
sharing, the absolute-majorant evaluation and the resulting fully
materialized equations. The basic Pell theorems and CRT are classical.

The report makes no claim of optimal degree/variable/height tradeoff,
historical uniqueness, short consistency proofs, or small possible
refutation witnesses. It supplies a fixed equation, an ordinary proof
of its exact target semantics, and reproducible finite artifacts.

\section{What this changes for the next iteration}

The main structural gain is a compatible common certificate for an entire
finite family of powers. The construction uses two global Pell norms and
two local norms per seed exponent. The theorem is reusable with a different
logical source; its proof states how all shared witnesses can be chosen.

\clearpage
\begin{thebibliography}{9}
\raggedright
\bibitem{repo}
Vladimir Reshetnikov and ProveIt contributors,
\emph{ProveIt}, inspected commit
\texttt{b594a5e988e986f695923d068ef3d33ab2ada16f}.
\url{https://github.com/VladimirReshetnikov/ProveIt/tree/b594a5e988e986f695923d068ef3d33ab2ada16f/Computability/DiophantineConsistency}.

\bibitem{latest}
ProveIt, \emph{Later arrivals and current size reductions},
\File{Computability/DiophantineConsistency/PROPOSAL_REVIEW_LATEST.md},
at the inspected commit. Includes the direct C222 and shared CRT
comparison points and the ordinary/formal proof boundaries.

\bibitem{size}
ProveIt, \emph{Active equation-size research},
\File{Computability/DiophantineConsistency/SIZE_REDUCTION.md},
at the inspected commit. Its table distinguishes domains, degree,
operation count and the native-coefficient quartic alternatives.

\bibitem{projection}
\emph{ConZFC Direct Projection 7617}, incoming archive
\File{docs/incoming/ConZFC_Direct_Projection_7617.zip},
at the inspected commit. Its literal rule graph, source projection proof,
arithmetic interfaces and logical pins are retained or explicitly
transformed in this package.

\bibitem{prime}
ProveIt, \emph{Smaller direct quartics from a prime-power CRT lift},
\File{candidates/shared-crt-prime-lift/PROOF.md},
under the inspected DiophantineConsistency directory.

\bibitem{calculus}
ProveIt, \File{KCode222.lean}, \File{KEndpoint.lean} and
\File{ConstructibleChoice.lean}, logical source pin
\texttt{9a33342a03d19cf5b061b2083c6b5d3f27f752ed}.
The full filenames, URLs and SHA-256 values are in
\File{provenance/LOGICAL_SOURCE_PINS.json}.

\bibitem{pell}
Mario Carneiro and the Mathlib contributors,
\emph{Pell's equation and Matiyasevic's theorem},
\File{Mathlib/NumberTheory/PellMatiyasevic.lean}.
\url{https://leanprover-community.github.io/mathlib4_docs/Mathlib/NumberTheory/PellMatiyasevic.html}.
The exact source blob used is copied in \File{provenance/}, with its
Apache 2.0 license.
\end{thebibliography}
\end{document}
